\documentclass[border=2pt]{standalone}
\usepackage[UTF8,fontset=none]{ctex}
\usepackage{tikz}
\usetikzlibrary{arrows.meta}
\setCJKmainfont[Path=../../../../assets/fonts/,BoldFont=LXGWWenKai-Medium.ttf]{LXGWWenKai-Regular.ttf}
\definecolor{navy}{HTML}{003F88}
\definecolor{sky}{HTML}{EAF2F9}
\definecolor{edge}{HTML}{6885A2}
\definecolor{teal}{HTML}{23685F}
\definecolor{mint}{HTML}{EAF4EF}
\definecolor{ink}{HTML}{20374B}
\tikzset{
  frame/.style={draw=edge,fill=sky,rounded corners=2pt,line width=.55pt},
  fixed/.style={draw=teal,fill=mint,rounded corners=2pt,line width=.55pt},
  link/.style={-{Stealth[length=2mm]},draw=edge,line width=.9pt},
  flowbox/.style={frame,align=center,minimum height=12mm,text width=31mm,inner sep=2pt},
  legendtext/.style={text=ink,anchor=west,font=\fontsize{9.5}{12}\selectfont},
  captionlabel/.style={text=navy,anchor=west,font=\bfseries\fontsize{11}{14}\selectfont}
}

\begin{document}
\begin{tikzpicture}[x=1cm,y=1cm,text=ink]
  % Two independent report paths. Lines only connect stages traversed by each path.
  \newcommand{\stage}[5]{
    \draw[#4] (#1,#2) rectangle ++(4.45,-1.65);
    \node[captionlabel] at (#1+.12,#2-.42) {#3};
    \node[legendtext,align=left] at (#1+.12,#2-1.05) {#5};}
  \node[captionlabel] at (0,0) {固定简报};
  \node[legendtext] at (2.05,0) {用已登记的资料验证三项指标};
  \stage{0}{-.35}{选择模板与资料}{fixed}{网页或命令行发起\\也可在命令行排除某来源}
  \stage{5.15}{-.35}{逐章查询}{fixed}{LangGraph 编排\\指标编号查固定资料}
  \stage{10.3}{-.35}{核对简报}{fixed}{命中、缺项均进入草稿\\人工确认后导出定稿}
  \draw[link] (4.45,-1.2)--(5.15,-1.2);
  \draw[link] (9.6,-1.2)--(10.3,-1.2);
  \node[captionlabel] at (0,-2.52) {自定义研究};
  \node[legendtext] at (2.05,-2.52) {研究员在本地网页指定题目、篇幅与资料范围};
  \stage{0}{-2.87}{拟定研究范围}{frame}{讨论提纲或直接提交\\按任务保存题目与选择}
  \stage{5.15}{-2.87}{收集资料并按章写作}{frame}{指定来源的原件或数据集\\提取证据并编号}
  \stage{10.3}{-2.87}{逐句审阅与修订}{frame}{查看原句、引用和原文\\草稿可润色，确认后导出}
  \draw[link] (4.45,-3.72)--(5.15,-3.72);
  \draw[link] (9.6,-3.72)--(10.3,-3.72);
\end{tikzpicture}
\end{document}
